
\item
Pour toute extension \(K\) de \(k\), le noyau de \(F(k)\to F(K)\) est un groupe de torsion.
\item Pour toute extension pure \(K\) de \(k\), l'homomorphisme \(F(k)\to F(K)\) est injectif. (NB. Noter qu'on ne suppose pas nécessairement \(k\) infini).
\end{enumerate}

Dans le cas (i), on est ramené à prouver que tout élément \(x\) dans le noyau de \(F(k)\to F(A)\), où \(A\) est une algèbre intègre de type fini sur \(k\), est un élément de torsion. Utilisant un \(k\)-homomorphisme \(A\to k'\), où \(k'\) est une \(k\)-algèbre finie, on est ramené au cas où \(A=k'\) est une extension finie de \(k\). Alors on a \([k':k]x=0\) en vertu de la propriété postulée plus haut pour l'homomorphisme norme. Cela prouve que \(x\) est un élément de torsion. On aura même \(x=0\) dans le cas où le pgcd des degrés des extensions finies \(k'\) de \(k\) quotients de \(A\) est égal à 1. Or cela est le cas si le corps des fractions \(K\) de \(A\) est une extension pure de \(k\), comme on constate aisément (seul le cas où \(k\) est fini demandant une nouvelle vérification). Cette remarque donne alors la conclusion de (ii).

\textbf{7.18.4.} Supposons que pour tout ouvert affine non vide \(U\) de l'espace affine \(\Spec k[T_1,\ldots,T_n]\), si \(A\) est l'anneau de \(U\), l'homomorphisme \(F(k)\to F(A)\) est surjectif (comparer IX 1.6). Alors, utilisant la commutation de \(F\) aux limites inductives filtrantes, on trouve que pour toute extension pure \(K\) de \(k\), \(F(k)\to F(K)\) est surjectif. Si on suppose de plus \(k\) infini, ou qu'on peut définir des homomorphismes norme comme dans 7.18.3, on en conclut que cet homomorphisme est même bijectif.

\textbf{7.18.5.} À titre d'exemple, considérons un schéma \(X\) de type fini sur \(k\), et posons
\[
  F(A)=K_\bullet(X\otimes_k A),
\]
